\documentclass[10pt,a4paper]{article}

\usepackage[T1]{fontenc}
\usepackage[utf8]{inputenc}
\usepackage[french]{babel}
\usepackage{helvet}
\renewcommand{\familydefault}{\sfdefault}
\usepackage[a4paper,left=10mm,right=10mm,top=9mm,bottom=8mm]{geometry}
\usepackage[table]{xcolor}
\usepackage{tabularx,booktabs}
\usepackage{enumitem}
\usepackage[most]{tcolorbox}
\usepackage{microtype}
\usepackage{hyperref}           
\hypersetup{colorlinks=true, urlcolor=orange, linkcolor=orange,
            pdftitle={Détection d'intrusion réseau explicable}, pdfauthor={Farah Grissa}}
\newcommand{\appli}{\href{https://ids-xai-classification.streamlit.app}{ids-xai-classification.streamlit.app}}
\newcommand{\api}{\href{https://ids-xai-api.onrender.com/docs}{ids-xai-api.onrender.com/docs}}
\newcommand{\github}{\href{https://github.com/Farahgrissa/ids-xai-classification}{github.com/Farahgrissa/ids-xai-classification}}

% --- Couleurs ---------------------------------------------------------------
\definecolor{orange}{HTML}{D9480F}
\definecolor{rouge}{HTML}{B42318}
\definecolor{encre}{HTML}{1B2430}
\definecolor{gris}{HTML}{5B6675}
\definecolor{trait}{HTML}{D9DEE5}
\definecolor{fondkpi}{HTML}{FFF4EC}
\definecolor{fondrouge}{HTML}{FDF3F2}
\definecolor{fondligne}{HTML}{FFE8D9}

\color{encre}
\pagestyle{empty}
\setlength{\parindent}{0pt}
\setlength{\parskip}{2pt}
\frenchsetup{StandardLists=true}  

% --- Titres de section ----------------------------------------------------------
\newcommand{\rubrique}[1]{%
  \vspace{4pt}%
  {\color{orange}\bfseries\small\MakeUppercase{#1}}\par\vspace{-3pt}%
  {\color{trait}\rule{\linewidth}{0.5pt}}\par\vspace{1pt}}

% --- Listes compactes -------------------------------------------------------------
\setlist[itemize]{leftmargin=3.2mm, itemsep=1.2pt, topsep=1pt, parsep=0pt,
                  label={\color{orange}\scriptsize$\bullet$}}

% --- Indicateurs clés ---------------------------------------------------------------
\newcommand{\kpi}[2]{%
  \begin{tcolorbox}[colback=fondkpi, colframe=fondkpi, arc=1.2mm, boxrule=0pt,
                    left=1.6mm, right=1.6mm, top=1.4mm, bottom=1.2mm, width=\linewidth,
                    height=15.5mm, valign=top]
    \raggedright
    {\color{orange}\bfseries\Large #1}\par\vspace{2pt}
    {\color{gris}\scriptsize #2}
  \end{tcolorbox}}

% --- Encadrés -----------------------------------------------------------------------
\newtcolorbox{alerte}{colback=fondrouge, colframe=rouge, boxrule=0pt, leftrule=0.9mm,
  arc=0pt, left=2mm, right=2mm, top=1.2mm, bottom=1.2mm, fontupper=\fontsize{8.7}{11.3}\selectfont}
\newtcolorbox{conseil}{colback=fondkpi, colframe=orange, boxrule=0pt, leftrule=0.9mm,
  arc=0pt, left=1.2mm, right=2mm, top=1mm, bottom=1mm, fontupper=\fontsize{8.7}{11.3}\selectfont}

\begin{document}
\fontsize{8.7}{11.3}\selectfont

% =================================================================================
% En-tête
% =================================================================================
{\color{orange}\bfseries\fontsize{19}{22}\selectfont Détection d'intrusion réseau explicable}\par
\vspace{2pt}
{\color{gris}\small Classification supervisée et explicable de flux réseau (CSE-CIC-IDS2018, infrastructure AWS) ---
Projet Challenge 1, Data Mining, ENSI 2026-2027 --- Farah Grissa}\par
\vspace{1pt}
{\color{gris}\scriptsize \textbf{\color{encre}Code :} \github
\quad·\quad \textbf{\color{encre}Application :} \appli
\quad·\quad \textbf{\color{encre}API :} \api}\par
\vspace{2pt}
{\color{orange}\rule{\linewidth}{1.4pt}}\par
\vspace{3pt}

% =================================================================================
% Indicateurs clés
% =================================================================================
\noindent
\begin{minipage}[t]{0.193\linewidth}\kpi{37\,395}{flux analysés après nettoyage (27\,\% d'attaques)}\end{minipage}\hfill
\begin{minipage}[t]{0.193\linewidth}\kpi{0,994}{PR-AUC de XGBoost sur un jeu de test jamais vu}\end{minipage}\hfill
\begin{minipage}[t]{0.193\linewidth}\kpi{98\,\%}{des attaques détectées (rappel, seuil 0,5)}\end{minipage}\hfill
\begin{minipage}[t]{0.193\linewidth}\kpi{$\approx$\,1\,\%}{du trafic normal signalé à tort}\end{minipage}\hfill
\begin{minipage}[t]{0.193\linewidth}\kpi{4 ms}{pour classer 1\,000 flux : compatible temps réel}\end{minipage}

\vspace{2pt}

% =================================================================================
% Colonne gauche
% =================================================================================
\noindent
\begin{minipage}[t]{0.485\linewidth}

\rubrique{Problème métier}
Un centre opérationnel de sécurité (SOC) reçoit des millions de flux réseau par jour. Il lui faut un détecteur
qui \textbf{repère les attaques sans noyer les analystes sous les fausses alertes}, et qui
\textbf{justifie chaque alerte} pour qu'elle soit vérifiée en quelques secondes.

\rubrique{Données}
\begin{itemize}
  \item \textbf{CSE-CIC-IDS2018}, trafic capturé sur AWS, téléchargé automatiquement depuis le bucket S3 public.
  \item 3 familles d'attaques : \textbf{force brute} (FTP, SSH), \textbf{déni de service} (Hulk, SlowHTTPTest),
        \textbf{botnet}. Cible binaire : normal / attaque.
  \item \textbf{Qualité :} 21\,\% de doublons supprimés ; étiquettes incohérentes détectées (flux identiques
        étiquetés FTP et SlowHTTPTest).
  \item \textbf{Raccourcis exclus :} heure, port de destination et variables d'« empreinte système »
        (taille de fenêtre TCP), qui identifient la machine plutôt que le comportement.
\end{itemize}

\rubrique{Démarche et meilleur modèle}
Cinq modèles comparés en validation croisée stratifiée (5 plis), hyperparamètres optimisés (GridSearch,
RandomizedSearch), critère \textbf{PR-AUC} adapté au déséquilibre des classes.

\vspace{2pt}
{\fontsize{8}{10}\selectfont
\renewcommand{\arraystretch}{1.2}
\begin{tabularx}{\linewidth}{@{}Xrrrr@{}}
  \toprule
  \textcolor{gris}{\textbf{Modèle (jeu de test)}} & \textcolor{gris}{\textbf{PR-AUC}} &
  \textcolor{gris}{\textbf{Rappel}} & \textcolor{gris}{\textbf{Précision}} & \textcolor{gris}{\textbf{ms / 1\,000}} \\
  \midrule
  \rowcolor{fondligne} \textbf{XGBoost (retenu)} & \textbf{0,994} & \textbf{0,982} & \textbf{0,973} & \textbf{4} \\
  Random Forest         & 0,993 & 0,981 & 0,976 & 11 \\
  Arbre de décision     & 0,985 & 0,979 & 0,968 & 3 \\
  KNN                   & 0,977 & 0,978 & 0,970 & 53 \\
  Régression logistique & 0,898 & 0,960 & 0,815 & 4 \\
  \bottomrule
\end{tabularx}}

\vspace{3pt}
XGBoost égale Random Forest en performance, prédit 2 à 3 fois plus vite et ne surapprend pas
(écart validation / test $<$ 0,003, donc aucune fuite de données).

\rubrique{Résultats en conditions réelles}
Sur le test, XGBoost manque environ 37 attaques sur 2\,043, principalement des DoS Hulk dont le serveur
saturé ne répond plus (flux vides, proches de connexions normales interrompues).

\begin{alerte}
\textbf{Effet du taux de base.} Le test contient 27\,\% d'attaques, un vrai réseau moins de 1\,\%.
À taux de fausses alertes égal, la précision passe alors de 97\,\% à \textbf{environ 50\,\%} :
une alerte sur deux serait fausse. Relever le seuil de 0,5 à 0,9 n'apporte que quelques points
de précision et double le nombre d'attaques manquées.
\end{alerte}

\end{minipage}\hfill
% =================================================================================
% Colonne droite
% =================================================================================
\begin{minipage}[t]{0.485\linewidth}

\rubrique{Pourquoi le modèle décide (SHAP, LIME)}
\begin{itemize}
  \item \textbf{Le rythme d'abord} (1re variable SHAP) : un intervalle de 0,1 à 50 ms entre paquets envoyés
        pousse vers l'attaque, au-delà de 0,1 s vers le normal. Effet en cloche, hors de portée d'un modèle linéaire.
  \item \textbf{Puis le contenu et TCP} : aucun octet échangé (requête sans réponse), connexion coupée par RST
        (+2 à +6 log-odds, signal le plus fort pour le botnet).
  \item \textbf{Une alerte expliquée} : Hulk à 98\,\% = deux paquets à 8,4 ms d'écart, aucun octet renvoyé.
        Une phrase de synthèse est générée pour chaque flux (application et API \texttt{/explain}).
  \item \textbf{Fausse alerte comprise} : un flux normal court et sans réponse (33 ms) a exactement le profil Hulk.
  \item \textbf{SHAP et LIME cohérents}, et avec l'intuition métier, si LIME est réglé localement : 2,8 variables
        communes sur 5, même sens. Avec son noyau par défaut, LIME n'est pas fidèle (R² 0,13).
\end{itemize}

\rubrique{Analyse critique}
\begin{itemize}
  \item \textbf{Variables sensibles et équité.} Aucune donnée personnelle (adresses IP exclues). L'équivalent est
        l'\textbf{empreinte système} : remise dans le modèle, elle devient la 1re raison SHAP (36\,\% de l'importance)
        et une fenêtre TCP de 225 donne 100\,\% d'attaques. Le modèle jugerait la machine, pas son comportement :
        elle est exclue (PR-AUC 0,9997 $\rightarrow$ 0,9936). Biais restant : l'UDP est appris comme normal.
  \item \textbf{Performance ou interprétabilité.} La régression logistique perd 10 points de PR-AUC, produit
        \textbf{9 fois plus de fausses alertes}, et ses contributions SHAP se compensent ($\pm$50 log-odds) :
        elle n'est pas plus lisible. Recommandation : \textbf{XGBoost avec une explication SHAP par alerte}.
\end{itemize}

\rubrique{Actions pour le décideur}
\begin{conseil}
\begin{itemize}
  \item \textbf{Seuil autour de 0,8}, ajusté à la capacité de traitement de l'équipe.
  \item \textbf{Afficher la phrase et les contributions SHAP} de chaque alerte pour vérifier en quelques secondes.
  \item \textbf{Agréger par source et par fenêtre de temps} : un flux court sans réponse isolé est une fausse
        alerte probable, une rafale une attaque.
  \item \textbf{Enrichir l'entraînement} (attaques UDP, FTP), recalibrer sur le trafic réel, surveiller la dérive
        des variables de rythme.
\end{itemize}
\end{conseil}

\rubrique{Limites}
\begin{itemize}
  \item Données de laboratoire et taux d'attaques irréaliste : la précision mesurée est optimiste.
  \item Classification flux par flux : la répétition, signature de la force brute, n'est pas exploitée.
  \item SHAP explique le modèle, pas la cause réelle ; LIME dépend du voisinage choisi.
\end{itemize}

\end{minipage}

% =================================================================================
% Pied de page
% =================================================================================
\vfill
{\color{trait}\rule{\linewidth}{0.5pt}}\par
{\color{gris}\scriptsize Source : CSE-CIC-IDS2018 (Sharafaldin et al., ICISSP 2018) -- Registry of Open Data on AWS
\hfill \github}

\end{document}